\chapter{Il sistema nervoso centrale (2)}

% ---------------------------------------------------------------------------
\section{La circolazione del sangue nel cervello}

Il cervello rappresenta il 2\% del peso corporeo totale, ma riceve il 15--20\% di
tutto il sangue in circolazione: il flusso deve restare sempre costante, perché i
neuroni sono cellule molto esigenti per il rifornimento di nutrienti ed energia.

\rubrica{Cosa porta e cosa rimuove il sangue}
\begin{itemize}
  \item il sangue \textbf{porta} al cervello: $O_2$, carboidrati, amminoacidi,
    grassi, ormoni, vitamine;
  \item il sangue \textbf{rimuove} dal cervello: biossido di carbonio, ammoniaca,
    lattato, ormoni.
\end{itemize}

\rubrica{Consumo metabolico}
\begin{itemize}
  \item l'encefalo è un organo ad alto consumo metabolico, con limitate risorse
    energetiche di glucosio e $O_2$: la vascolarizzazione arteriosa deve quindi
    essere costante e metabolismo-dipendente;
  \item al cervello arrivano circa 750\,cc di sangue al minuto (50\,cc/100\,g,
    contro i 100\,cc/100\,g del cuore);
  \item il cervello consuma circa 50\,cc di $O_2$/min (quasi il 15\% del totale)
    e 80\,mg/min di glucosio.
\end{itemize}

% ---------------------------------------------------------------------------
\section{Vascolarizzazione arteriosa dell'encefalo}

Alla base della scatola cranica è presente un sistema di anastomosi arteriose a
pieno canale, dato dalla confluenza di 3 arterie principali:
\begin{itemize}
  \item \textbf{l'arteria basilare}: formata dalla confluenza delle arterie
    vertebrali destra e sinistra (prime collaterali della succlavia);
  \item \textbf{le 2 arterie carotidi interne} (destra e sinistra).
\end{itemize}

\rubrica{Il poligono di Willis}
Può essere ricondotto idealmente a un eptagono, i cui lati sono:
\begin{itemize}
  \item anteriormente, le 2 arterie cerebrali anteriori (dx e sx), unite
    dall'arteria comunicante anteriore;
  \item posteriormente, le 2 arterie cerebrali posteriori (dx e sx);
  \item tra arteria cerebrale anteriore e posteriore di ogni lato, l'arteria
    comunicante posteriore.
\end{itemize}

\rubrica{Varianti anatomiche}
Il sistema arterioso può presentare varianti:
\begin{itemize}
  \item a livello di origine;
  \item di decorso (tortuosità);
  \item di calibro;
  \item del sifone carotideo.
\end{itemize}

\rubrica{Occlusione della circolazione cerebrale}
Quando la circolazione del sangue nel cervello è bloccata si parla di ictus,
trombosi, embolo o stenosi.

% ---------------------------------------------------------------------------
\section{I seni venosi}

I seni venosi sono canali a pareti rigide situati nello spessore della dura
madre, rivestiti da endotelio.

\rubrica{Caratteristiche}
\begin{itemize}
  \item volume più considerevole rispetto alle arterie;
  \item i rami decorrono di preferenza alla superficie libera delle
    circonvoluzioni;
  \item pareti molto sottili, prive di fibre muscolari;
  \item assenza di valvole;
  \item molteplicità di anastomosi;
  \item sono tributari, direttamente o indirettamente, delle vene giugulari
    interne;
  \item ricevono sangue dalle vene cerebrali.
\end{itemize}

% ---------------------------------------------------------------------------
\section{I nuclei della base}

Localizzati nel telencefalo basale, costituiscono nel loro insieme il
\textbf{corpo striato}, suddiviso filogeneticamente in:
\begin{itemize}
  \item \textbf{striato} o neo-striato: nucleo caudato e putamen;
  \item \textbf{pallido} o paleo-striato: \textit{globus pallidus} laterale e
    mediale.
\end{itemize}

Rappresentano un importante sistema di selezione, avvio e arresto del
movimento volontario.

\rubrica{Circuiti}
Motorio, oculomotorio, cognitivo, limbico.

\rubrica{Funzioni}
\begin{itemize}
  \item controllo inconscio e integrazione del tono della muscolatura
    scheletrica;
  \item coordinazione degli schemi motori acquisiti;
  \item elaborazione, integrazione e trasmissione delle informazioni dalla
    corteccia cerebrale verso il talamo.
\end{itemize}

\rubrica{Patologie dei nuclei della base}
Sindromi rigido-acinetiche, sindromi iper-acinetiche, disturbi del movimento.

% ---------------------------------------------------------------------------
\section{Il sistema limbico}

Comprende nuclei e fasci che si trovano lungo il confine tra cervello e
diencefalo: rappresenta un raggruppamento più funzionale che anatomico, legato
al comportamento dell'individuo.

\rubrica{Componenti}
\begin{itemize}
  \item \textbf{funzioni}: elaborazione dei ricordi, creazione degli stati
    emozionali, della condotta e del comportamento;
  \item \textbf{componenti cerebrali}:
    \begin{itemize}
      \item aree corticali -- lobo limbico (giro del cingolo, giro
        paraippocampale, giro dentato);
      \item nuclei -- ippocampo e corpo amigdaloideo;
      \item tratti -- fornice;
    \end{itemize}
  \item \textbf{componenti diencefaliche}:
    \begin{itemize}
      \item talamo -- gruppo nucleare anteriore;
      \item ipotalamo -- centri delle emozioni, dell'appetito (fame e sete) e
        del relativo comportamento;
    \end{itemize}
  \item \textbf{altre componenti}: formazione reticolare, rete di nuclei
    cerebrali interconnessi in tutto il tronco encefalico.
\end{itemize}

% ---------------------------------------------------------------------------
\section{Il talamo}

Rappresenta il punto di trasmissione finale per le informazioni sensitive e
coordina le attività motorie somatiche volontarie e involontarie. Contiene
nuclei che rappresentano centri di controllo delle informazioni sensitive e
motorie.

\begin{itemize}
  \item forma simile a un uovo: è il ``relè'' del cervello;
  \item struttura posta tra tronco encefalico ed emisferi cerebrali;
  \item quasi tutte le informazioni sensoriali passano da qui prima di essere
    inviate alla corteccia;
  \item quasi tutti i nuclei talamici sono connessi con la corteccia.
\end{itemize}

% ---------------------------------------------------------------------------
\section{L'ipotalamo}

Controlla numerose funzioni autonome, tra cui l'omeostasi e il controllo
ipofisario: è il centro di collegamento tra sistema nervoso e sistema
endocrino.

\rubrica{Funzioni}
\begin{enumerate}
  \item \textbf{secrezione ormonale}: secerne ossitocina e ormone
    antidiuretico;
  \item \textbf{effetti sul sistema nervoso autonomo}: è uno dei principali
    centri di integrazione del SNA; invia fibre discendenti ai nuclei del
    tronco encefalico che influenzano ritmo cardiaco, pressione arteriosa,
    diametro delle pupille, motilità gastrica, secrezione gastrica e altre
    funzioni;
  \item \textbf{termoregolazione}: i neuroni del nucleo preottico monitorano
    la temperatura corporea; quando questa si allontana dalla norma vengono
    attivati meccanismi per alzarla o abbassarla (sudorazione,
    vasodilatazione o vasocostrizione cutanea);
  \item \textbf{controllo dei ritmi circadiani}: la regione caudale
    ipotalamica fa parte della formazione reticolare, che contiene nuclei
    coinvolti nella regolazione del ritmo sonno-veglia;
  \item \textbf{assunzione di cibo e acqua}: i neuroni ipotalamici
    (osmocettori) controllano l'osmolarità del sangue e, in caso di
    disidratazione, innescano comportamenti finalizzati alla ricerca di
    acqua o cibo.
\end{enumerate}

% ---------------------------------------------------------------------------
\section{Il cervelletto}

Occupa la fossa cranica posteriore e rappresenta la seconda porzione più
estesa del cervello in toto: un ``encefalo mignon'', costituito da 2 emisferi
cerebellari connessi da uno stretto ponte chiamato \textbf{verme}.

\begin{itemize}
  \item presenta una corteccia esterna di sostanza grigia e una sostanza
    bianca interna, con morfologia simile a una felce, chiamata
    \textbf{arbor vitae};
  \item è collocato posteriormente al tronco dell'encefalo, a cui è collegato
    tramite i peduncoli cerebellari;
  \item forma il tetto del IV ventricolo.
\end{itemize}

\rubrica{Le tre regioni}
\begin{itemize}
  \item \textbf{corteccia}: sostanza grigia (più superficiale);
  \item \textbf{sostanza bianca} interna (più in profondità);
  \item \textbf{nuclei profondi}: materia grigia in profondità.
\end{itemize}

Contiene molte cellule di Purkinje, dalle complesse ramificazioni.

\rubrica{Dimensioni}
Forma grossolanamente ovoidale, con grande asse trasversale di 9\,cm,
diametro antero-posteriore di 6\,cm e verticale di 5\,cm; peso medio 140\,g.
È suddiviso nei lobi flocculonodulare, posteriore e anteriore.

\rubrica{Le cellule del Purkinje}
\begin{itemize}
  \item grossi neuroni il cui corpo oblungo è situato nell'omonimo strato;
  \item disposte a intervalli piuttosto regolari;
  \item possiedono un albero dendritico fittamente arborizzato, che origina
    da uno o due dendriti primari e si porta fino allo strato molecolare,
    dove contrae migliaia di sinapsi tramite numerose spine dendritiche con
    le fibre parallele delle cellule dei granuli;
  \item l'albero dendritico è caratteristicamente limitato a una sottile
    lamina di tessuto perpendicolare all'asse maggiore della lamella.
\end{itemize}

\rubrica{Ruolo funzionale}
Attaccato alla parte posteriore dell'encefalo, simile a un cavolfiore e
grande come un pugno; caratterizzato da altissima densità cellulare, è il
10\% del volume cerebrale ma contiene più o meno metà dei neuroni del SNC.
È coinvolto in \textit{movimento}, \textit{equilibrio} e \textit{postura};
oggi si ritiene che partecipi anche a memoria, umore, linguaggio, attenzione
e funzioni sensitive/emotive.

% ---------------------------------------------------------------------------
\section{Il midollo spinale}

Porzione extracranica del SNC, collocata all'interno del canale vertebrale.

\begin{itemize}
  \item comincia dal foro occipitale, come prosecuzione del bulbo, e continua
    in senso caudale fino al cono midollare a livello lombare;
  \item la sua estensione fibrosa, detta \textbf{filum terminale}, si
    prolunga fino al coccige;
  \item lunghezza media nell'adulto: 46\,cm (larghezza 1--1,5\,cm), in
    stretto rapporto con la lunghezza del tronco e variabile secondo sesso ed
    etnia;
  \item il livello di terminazione è variabile: nella maggior parte dei casi
    l'apice inferiore si trova nel disco intervertebrale tra L1 e L2, a
    livello del piano transpilorico, ma può terminare anche tra L2 e L3, o
    più raramente tra T12 e L1.
\end{itemize}

\rubrica{Organizzazione}
È formato da sostanza bianca esterna (prolungamenti dei neuroni) e sostanza
grigia interna (corpi dei neuroni), ed è diviso in quattro regioni --
cervicale, toracica, lombare e sacrale -- ciascuna formata da più segmenti.

\rubrica{Meningi spinali}
La dura madre riveste il midollo spinale formando un sacchetto di protezione,
irrorato al suo interno dal liquido cerebrospinale.

\rubrica{Legamenti denticolati}
Coppie di estensioni della pia madre spinale che collegano pia madre e
aracnoide spinali alla dura madre; impediscono i movimenti di scivolamento
verso il basso del midollo spinale.

\rubrica{Sostanza grigia}
A forma di H:
\begin{itemize}
  \item \textbf{commissura grigia}: contiene il canale centrale;
  \item \textbf{corna anteriori}: contengono i corpi cellulari dei
    motoneuroni;
  \item \textbf{corna posteriori}: consistono di interneuroni;
  \item la sostanza grigia è suddivisa tra regioni viscerali e somatiche.
\end{itemize}

\rubrica{Sostanza bianca}
Formata da assoni mielinici e amielinici, organizzati in tre tipi di fibre:
ascendenti, discendenti e commissurali.
